\documentclass[11pt]{article}
\usepackage[T1]{fontenc}
\usepackage{textcomp}
\usepackage[margin=0.75in]{geometry}
\usepackage{amsmath,amssymb,amsthm}
\usepackage{array,booktabs}
\usepackage{hyperref}
\setlength{\parindent}{0pt}
\newtheorem{theorem}{Theorem}
\theoremstyle{definition}
\newtheorem{definition}{Definition}
\title{Flow Proof Artifact: Nat-mul-distrib-right}
\author{Generated by \texttt{flow doc proof}}
\date{Flow Proof Book}
\begin{document}
\maketitle
Multiplication distributes over addition on the left.\\[0.5em]
\textbf{Source.} peano — https://en.wikipedia.org/wiki/Distributive\_property\\[0.5em]
\bigskip
\noindent{\large \textbf{Derived fact 1.} \textit{(a + b) * c = a * c + b * c} \hfill the natural numbers \cdot multiplication \cdot distributes over addition on the left}\par
\smallskip\noindent\textit{Coordinate: the natural numbers \cdot multiplication \cdot distributes over addition on the left}\par
\smallskip\noindent\textit{Source: peano/induction — Gries \& Schneider, Ch. 3}\par
\smallskip\noindent\textit{Built on: distributes over addition on the right, for multiplication on the natural numbers, order does not matter, for multiplication on the natural numbers}\par
\medskip
\noindent\textbf{Goal.} (a + b) * c = a * c + b * c\par
\noindent$\displaystyle \forall a \in \mathbb{N} \forall b \in \mathbb{N} \forall c \in \mathbb{N}\quad (a + b) * c = a \cdot c + b \cdot c$\par
\medskip
\noindent
\begin{tabular}{@{} >{\raggedright\arraybackslash}p{0.06\textwidth} >{\raggedright\arraybackslash}p{0.47\textwidth} >{\raggedleft\arraybackslash}p{0.41\textwidth} @{}}
\textbf{\#} & \textbf{Proof} & \textbf{Mathematics} \\
\midrule
\textbf{1.} & We prove that distributes over addition on the left for multiplication on the natural numbers. & \\[0.45em]
\textbf{2.} & We invoke the derived fact governing multiplication on the natural numbers: distributes over addition on the right, for multiplication on the natural numbers (instantiated for c, a, b). & \\[0.45em]
\textbf{3.} & We invoke the derived fact governing multiplication on the natural numbers: order does not matter, for multiplication on the natural numbers (instantiated for a, c). & \\[0.45em]
\textbf{4.} & We invoke the derived fact governing multiplication on the natural numbers: order does not matter, for multiplication on the natural numbers (instantiated for b, c). & \\[0.45em]
\textbf{5.} & We invoke the derived fact governing multiplication on the natural numbers: order does not matter, for multiplication on the natural numbers (instantiated for a + b, c). & \\[0.45em]
\textbf{6.} & We invoke the derived fact governing multiplication on the natural numbers: order does not matter, for multiplication on the natural numbers (instantiated for c, a + b). & \\[0.45em]
\textbf{7.} & From step 2, step 3, step 4, step 5, and step 6, this implies (a plus b) times c equals a times c plus b times c. Hence proven. & $\displaystyle (a + b) * c = a \cdot c + b \cdot c$ \\[0.45em]
\bottomrule
\end{tabular}
\label{thm:Nat-multiplication-distributes_over_addition_on_the_left}
\par\medskip\hrule
\end{document}
